\chapter{Adenoidectomy}
\index{Adenoidectomy}

\section*{Definition and Overview}
An adenoidectomy is an operation to remove the adenoids, a pad of immune tissue at the back of the nose. It is performed most often in children when enlarged adenoids block the airway or the passage between the nose and the middle ear, causing noisy breathing, sleep problems, glue ear, or repeated ear infections.

The operation is performed under a general anaesthetic through the mouth, takes about 20 to 30 minutes, and most children go home the same day. It may be done alone or together with the insertion of grommets, or with removal of the tonsils (adeno-tonsillectomy).

\section*{Epidemiology}
Adenoidectomy is one of the more common operations in children, usually performed between the ages of 1 and 7 years, when the adenoids are naturally at their largest. It is performed less often in older children and adults, because the adenoids shrink during the teenage years.

Surgery rates are higher in children with persistent middle-ear disease, glue ear, or obstructive sleep apnoea, and vary between countries and health systems depending on referral practices and access to services.

\section*{Aetiology and Pathophysiology}
The adenoids are part of the body's immune defences and enlarge as they fight infection. In some children they remain large and do not shrink as expected, blocking the back of the nose. This makes mouth breathing necessary, disrupts sleep, and can block the opening of the eustachian tube, which drains fluid from the middle ear.

When the tube is blocked, fluid collects behind the eardrum (glue ear), causing hearing problems and making middle-ear infections more likely. Over time, persistent blockage can also contribute to speech delay and disturbed sleep with effects on growth and behaviour.

\section*{Clinical Features}
Children who need an adenoidectomy often have:

\begin{itemize}
    \item \textbf{Nasal blockage:} persistent mouth breathing and a blocked or stuffy nose
    \item \textbf{Snoring:} loud snoring and restless sleep
    \item \textbf{Breathing pauses:} pauses in breathing during sleep (obstructive sleep apnoea)
    \item \textbf{Voice change:} a nasal or muffled voice
    \item \textbf{Ear infections:} recurrent middle-ear infections
    \item \textbf{Glue ear:} fluid behind the eardrum, with reduced hearing and speech delay
    \item \textbf{Recurrent colds:} frequent colds or a constantly runny nose
\end{itemize}

\section*{Differential Diagnosis}
Symptoms of enlarged adenoids can be caused by other conditions:

\begin{itemize}
    \item Allergic rhinitis (hay fever)
    \item Enlarged tonsils
    \item Chronic sinusitis
    \item Nasal polyps or a deviated septum
    \item A foreign body in the nose (in young children)
    \item Obesity-related obstructive sleep apnoea
\end{itemize}

\section*{When to Seek Medical Care}
See a doctor if a child has persistent mouth breathing, heavy snoring, or concerns about hearing, speech, or frequent ear infections. After surgery, contact the surgical team if the child has any bleeding, cannot drink enough, develops a fever, increasing pain, or a foul smell from the mouth, or if breathing is difficult or very noisy.

\textbf{Contact your local emergency services for:}
\begin{itemize}
    \item Bright red bleeding from the mouth or nose
    \item Difficulty breathing or choking
    \item Unusual drowsiness or difficulty waking
\end{itemize}

\section*{Diagnosis and Investigations}
\begin{itemize}
    \item \textbf{History:} symptoms of nasal blockage, sleep disturbance, hearing problems, and ear infections
    \item \textbf{Examination:} looking into the nose and throat
    \item \textbf{Nasendoscopy:} a small flexible camera to see the adenoids directly
    \item \textbf{Lateral X-ray of the neck:} may show how large the adenoids are
    \item \textbf{Hearing tests:} to assess glue ear or hearing loss
    \item \textbf{Sleep study:} for children with suspected obstructive sleep apnoea
\end{itemize}

\section*{Management}
\begin{itemize}
    \item \textbf{Watchful waiting:} many children improve as the adenoids shrink with age
    \item \textbf{Nasal steroid sprays:} may reduce adenoid size and nasal symptoms
    \item \textbf{Treatment of allergies:} antihistamines and avoiding triggers where relevant
    \item \textbf{Surgery (adenoidectomy):} for persistent obstruction, sleep apnoea, glue ear, or repeated ear infections
    \item \textbf{Grommets:} sometimes inserted at the same time to drain the middle ear
    \item \textbf{Post-operative care:} pain relief, fluids, a soft diet, and rest for about a week
\end{itemize}

\section*{Complications}
\begin{itemize}
    \item Bleeding after surgery, which is uncommon but can require treatment
    \item Sore throat and ear pain lasting up to a week or two
    \item Voice changes; some children briefly sound nasal
    \item Velopharyngeal insufficiency, where food or fluid comes up the nose, usually temporary
    \item Infection, and rarely regrowth of adenoid tissue
    \item Complications of the general anaesthetic, which are rare
\end{itemize}

\section*{Prognosis and Outcomes}
Adenoidectomy is very effective for its main indications. Most children breathe more easily through the nose, snore less, and sleep better afterwards. Where glue ear has caused hearing loss, hearing and speech often improve, and children may have fewer ear infections.

Recovery is usually straightforward, with most children back to normal activity within about a week or two. As with any surgery, outcomes are best when the decision is made on clear evidence of persistent obstruction, infection, or hearing problems rather than on a single episode.

\section*{Prevention and Self-Care}
\begin{itemize}
    \item Treat allergies and nasal congestion early to reduce adenoid symptoms
    \item Avoid tobacco smoke around children, which irritates the airways
    \item Practise good hand hygiene to reduce the spread of infections
    \item Encourage a healthy diet to support the immune system
    \item Attend follow-up and hearing checks as arranged
    \item Seek advice early about snoring, mouth breathing, or hearing concerns
\end{itemize}

\section*{Sources and Further Reading}
The following organisations provide reliable information on adenoidectomy and enlarged adenoids.

\begin{itemize}
    \item NHS. \textit{Information on adenoids, glue ear, and adenoidectomy.} https://www.nhs.uk/conditions/
    \item Mayo Clinic. \textit{Guide to adenoidectomy in children.} https://www.mayoclinic.org/
    \item MedlinePlus. \textit{Health information on adenoid removal.} https://medlineplus.gov/
    \item NICE. \textit{Clinical guidance on otitis media, glue ear, and surgical management.} https://www.nice.org.uk/
    \item MSD Manual. \textit{Professional information on ear, nose, and throat disorders in children.} https://www.msdmanuals.com/
\end{itemize}
